\section{Estrategia de datos: además de la escala, la calidad, la estructura y las etapas de aprendizaje}

La sección anterior expuso cómo calcula el modelo; esta sección se pregunta de qué distribución aprende realmente el modelo. La estrategia de datos no se reduce a «recolectar más páginas web»: el source mix, el filtrado por calidad, el orden, la tasa de repetición y la proporción por etapas modifican el gradiente. La clase primero plantea el problema mediante una comparación con el aprendizaje humano y después muestra, con dos proyectos de investigación concretos, cómo convertir la intuición en experimentos.

\subsection{Por qué los datos son el backbone y no un depósito de combustible}

La sección anterior explicó cómo procesa un block los tokens; ahora hay que dirigir la atención a la distribución de muestras que genera el gradiente. Se dice que los datos son el backbone porque el modelo interioriza como comportamiento predeterminado los hechos, formatos y estrategias que aparecen con alta frecuencia, en lugar de consultarlos uno por uno como una base de datos al terminar el entrenamiento. Al leer este apartado, distinga primero entre la «cantidad» del corpus y la «densidad efectiva de supervisión»: un mismo token, si está repetido, es erróneo o no guarda relación con el objetivo, puede aportar a la capacidad una contribución marginal completamente distinta.

Para una distribución de datos $q(x)$, el objetivo de entrenamiento optimiza en realidad
\[
\mathbb{E}_{x\sim q}\big[\ell_{\theta}(x)\big].
\]
Cambiar $q$ cambia los conocimientos, el estilo, el idioma y los patrones de razonamiento que el modelo ve recompensados una y otra vez. Aumentar el volumen de datos solo amplía el número de muestras; si el ruido, la repetición o el sesgo de dominio crecen al mismo tiempo, los tokens adicionales pueden producir rendimientos decrecientes e incluso sobrescribir capacidades ya adquiridas.


\lecturefigure{slide-024.jpg}{Impacto de los datos en la IA: acoplamiento entre calidad, cómputo, escala y aplicaciones}{CS25 V5 official deck, p.~24}


\lecturefigure{slide-025.jpg}{Estrategias de datos más inteligentes: prioridad a la calidad, proporciones por etapa y selección activa}{CS25 V5 official deck, p.~25}

\begin{knowledgebox}{Lectura de la figura: cómo leer la cadena de impacto de los datos}
Primero siga la cadena causal de la slide 24, de data a model behavior, y después lea la lista de estrategias de la slide 25. El primer paso es comparar quality y diversity: la primera reduce el ruido; la segunda evita que la capacidad se estreche. El segundo paso es ver si la selection consiste en filtrar antes del entrenamiento, hacer sampling durante el entrenamiento o realizar active collection después del entrenamiento; las tres etapas tienen costos distintos. El tercer paso es comprobar si la metric coincide con la capacidad objetivo. La figura respalda que «las decisiones sobre los datos influyen en el modelo», pero de un diagrama de flujo no puede deducirse que cierto filtro sea eficaz en todas las scales.
\end{knowledgebox}

\begin{importantbox}{Los cinco controles de la estrategia de datos}
\begin{enumerate}
  \item \textbf{Quality}: si las muestras son correctas, coherentes y de alta densidad informativa;
  \item \textbf{Diversity}: si los idiomas, dominios, tareas y formas de expresión cubren la cola larga;
  \item \textbf{Structure}: formas secuenciales como diálogos, cuentos, código y demostraciones;
  \item \textbf{Scale}: número de tokens y número de repeticiones;
  \item \textbf{Schedule}: proporción en que aparecen los distintos datos en las fases tempranas y tardías del entrenamiento.
\end{enumerate}
Un mismo corpus puede dar lugar a modelos distintos con schedules distintos.
\end{importantbox}

\subsection{Humanos y LLM: las diferencias son hipótesis de investigación, no pruebas de semejanza con el cerebro}

La clase contrasta baby/ChatGPT y brain/neural network para plantear un conjunto de preguntas abiertas. Los humanos aprenden de forma continua, actúan con objetivos, reciben percepciones multimodales y obtienen retroalimentación en un entorno social; los LLM clásicos optimizan el next-token objective principalmente en un único preentrenamiento fuera de línea a gran escala. El valor de esta comparación es generar ejes experimentales, no afirmar que los mecanismos biológicos ya se comprenden.


\lecturefigure{slide-026.jpg}{Baby y ChatGPT, brain y neural network: comparación entre dos tipos de sistemas de aprendizaje}{CS25 V5 official deck, p.~26}


\lecturefigure{slide-027.jpg}{Diferencias entre el aprendizaje humano continuo, orientado a objetivos, multimodal y con inducción estructurada, y el entrenamiento en una sola etapa de los LLM}{CS25 V5 official deck, p.~27}

\teachervoice{00:12:38--00:14:46, Steven enumera el aprendizaje continuous, goal-directed, multimodal y hierarchical como posibles razones de que los humanos sean más eficientes, y al mismo tiempo usa repetidamente «potential» y «might». Las notas conservan estos puntos como hypotheses y no los presentan como conclusiones de la neurociencia.}

\begin{warningbox}{Tres trampas al comparar la eficiencia de datos}
La supervisión que reciben los humanos no se limita a caracteres contables; la visión, la acción, la recompensa, la interacción social y la retroalimentación del entorno a largo plazo difícilmente pueden convertirse en tokens. Además, las funciones objetivo del modelo y del humano son distintas. Por último, los límites de la capacidad de «saber hablar» no son equivalentes. Por ello, «los niños reciben muchos menos tokens» sirve para plantear una pregunta, pero no demuestra por sí solo que cierta fuente de datos sea una explicación suficiente.
\end{warningbox}


\lecturefigure{slide-028.jpg}{Cuentos y diálogos infantiles: datos muy estructurados, muy interactivos y orientados al aprendiz}{CS25 V5 official deck, p.~28}


\lecturefigure{slide-029.jpg}{De internet a ChatGPT: un corpus web masivo pero heterogéneo}{CS25 V5 official deck, p.~29}

\subsection{Por qué estudiar modelos pequeños y datos pequeños}

La comparación human/LLM anterior señala muchos factores posibles, pero si se validan directamente con el modelo más grande, el presupuesto solo alcanza para uno o dos puntos, y cualquier conclusión quedará entrelazada con los hiperparámetros y la aleatoriedad. El valor científico de los modelos pequeños está en descomponer un problema grande en curvas densas: repetir el entrenamiento, sustituir datos y observar la convergencia y los tipos de error. También son un objetivo de deployment, porque los modelos locales inciden en la privacidad, la latencia y la accesibilidad; sin embargo, estos dos valores deben separarse: que los modelos pequeños sean fáciles de estudiar no permite suponer que su ordenamiento de scaling coincida necesariamente con el de los modelos grandes.

Los experimentos a pequeña escala reducen el costo de reproducción, permiten a los investigadores ejecutar varios seeds, hacer ablations más densas y examinar la dinámica del entrenamiento; también favorecen el uso on-device, la privacidad y la investigación abierta. Pero los resultados de modelos pequeños no pueden extrapolarse directamente al trillion-token regime: hay que distinguir entre «descubrir mecanismos candidatos» y «demostrar beneficios a gran escala».


\lecturefigure{slide-030.jpg}{Eficiencia, interpretabilidad, apertura y valor científico de estudiar modelos pequeños y datos pequeños}{CS25 V5 official deck, p.~30}

\begin{knowledgebox}{Lectura de la figura: la investigación con modelos pequeños tiene dos cadenas de valor}
Una es el \textbf{scientific proxy}: múltiples seeds a bajo costo, ablations densas y observación de la learning curve; la otra es el \textbf{deployment target}: on-device, baja latencia, privacidad y accesibilidad. Al leer la slide 30, no mezcle ambas cadenas. Que un método tenga un mecanismo claro en un modelo de 100M no implica que su ordenamiento se mantenga en uno de 70B; que un modelo de 3B pueda ejecutarse localmente tampoco implica que sea adecuado para responder preguntas de alto riesgo. La conclusión correcta es que la pequeña escala ofrece mejor resolución experimental y establece una Pareto frontier independiente para el despliegue.
\end{knowledgebox}

\teachervoice{00:13:55--00:14:46, el docente vincula el valor de los modelos pequeños con la ejecución local, la controlabilidad, el acceso al código abierto y la investigación cognitiva. Experiencia práctica: el producto más valioso de los experimentos pequeños son tendencias verificables y casos de falla, no una simulación barata de un entrenamiento grande.}

\subsection{Resumen de la sección}

La estrategia de datos abarca a la vez calidad, estructura, diversidad, escala y schedule. La comparación con el aprendizaje humano aporta dimensiones comprobables, y los modelos pequeños ofrecen un campo experimental de bajo costo. Las dos investigaciones siguientes responden, respectivamente, si los corpus orientados a niños son más eficaces y si los datos de alta calidad deben concentrarse en la fase final del entrenamiento.
